%%%%%%%% ICML 2025 EXAMPLE LATEX SUBMISSION FILE %%%%%%%%%%%%%%%%%

\documentclass{article}

% Recommended, but optional, packages for figures and better typesetting:
\usepackage{microtype}
\usepackage{graphicx}
\usepackage{subfigure}
\usepackage{booktabs} % for professional tables

% hyperref makes hyperlinks in the resulting PDF.
\usepackage[colorlinks]{hyperref}

% Use the following line for the initial blind version submitted for review:
\usepackage{icml2026}

% If accepted, instead use the following line for the camera-ready submission:
%\usepackage[accepted]{icml2026}

% For theorems and such
\usepackage{amsmath}
\usepackage{amssymb}
\usepackage{mathtools}
\usepackage{amsthm}

% if you use cleveref..
\usepackage[capitalize,noabbrev]{cleveref}

%%%%%%%%%%%%%%%%%%%%%%%%%%%%%%%%
% THEOREMS
%%%%%%%%%%%%%%%%%%%%%%%%%%%%%%%%
\theoremstyle{plain}
\newtheorem{theorem}{Theorem}[section]
\newtheorem{proposition}[theorem]{Proposition}
\newtheorem{lemma}[theorem]{Lemma}
\newtheorem{corollary}[theorem]{Corollary}
\theoremstyle{definition}
\newtheorem{definition}[theorem]{Definition}
\newtheorem{assumption}[theorem]{Assumption}
\theoremstyle{remark}
\newtheorem{remark}[theorem]{Remark}

% Todonotes is useful during development; simply uncomment the next line
%    and comment out the line below the next line to turn off comments
%\usepackage[disable,textsize=tiny]{todonotes}
% \usepackage[textsize=tiny]{todonotes}

% Running title (short version)
\icmltitlerunning{Token-Level Subspace Routing}

\begin{document}

\twocolumn[
\icmltitle{Token-Level Subspace Routing:\\
          A New Paradigm for Transformer Hidden Representations}

% Equal contribution symbol
\icmlsetsymbol{equal}{*}

\begin{icmlauthorlist}
\icmlauthor{Chengkai Zhang}{equal,jhu}
\icmlauthor{Shujie Liu}{jhu}
\icmlauthor{Muyun Yang}{hit}  % 新增作者，倒数第二位
\icmlauthor{Xiaohua Liu}{equal,jhu,yyx}
\end{icmlauthorlist}

\icmlaffiliation{jhu}{Jianghan University, Wuhan, Hubei, China}
\icmlaffiliation{hit}{Harbin Institute of Technology, Harbin, Heilongjiang, China}  % 新增哈尔滨工业大学
\icmlaffiliation{yyx}{Yuanyu Xinqing (Xiong'an) Technology Co., Ltd., Xiong'an New Area, Hebei, China}

\icmlcorrespondingauthor{Xiaohua Liu}{xiaohualiu@jhun.edu.cn}


\icmlkeywords{Transformers, Subspace Routing, Sparse Computation, Representation Learning}


\vskip 0.3in
]

% This actually prints the affiliations + equal contribution footnote
\printAffiliationsAndNotice{\icmlEqualContribution}


\begin{abstract}
Transformers allocate the same $d$-dimensional hidden representation to every
token, regardless of its semantic complexity. This uniform capacity allocation
is a long-standing architectural assumption that conflicts with empirical
evidence: token activations and hidden features concentrate on highly
structured, low-dimensional subspaces. We identify \emph{token-level capacity
allocation} as a fundamental but previously unaddressed modeling problem.

We introduce \emph{Token-Level Subspace Routing} (TLSR), a new paradigm in
which each token dynamically activates only a small subset of learned hidden
subspaces. TLSR factorizes the hidden dimension, learns a token-dependent
routing mechanism, and induces structured block-sparse computations for both
FFN and attention projections—while keeping the Transformer architecture
unchanged. Unlike MoE or activation sparsity, TLSR redistributes existing
capacity at a finer granularity without introducing experts or unstructured
sparsity.

Empirically and theoretically, we show that token representations inhabit
multiple low-dimensional manifolds, and that TLSR aligns computation with this
geometry while enabling efficient block-sparse GPU kernels. Across large
language model benchmarks, TLSR matches or improves accuracy while reducing
FLOPs and wall-clock latency % [TLSR_RESULTS], establishing token-level subspace routing as a
promising and general direction for scaling efficient Transformers.
\end{abstract}



\section{Introduction}

Modern Transformer-based large language models (LLMs) allocate the same
$d$-dimensional hidden representation to every token at every layer.
This architectural choice---an implicit ``uniform capacity allocation
assumption''---has remained unchanged since the original Transformer design,
despite the fact that different tokens vary dramatically in semantic
complexity, contextual dependence, and intrinsic dimensionality.
Common function words, punctuation, and structural markers rarely require the
same representational capacity as content-bearing tokens, yet all tokens are
processed with identical full-dimensional linear projections.

Recent analyses provide growing evidence that Transformers do not use their
hidden dimensions uniformly.
A substantial portion of neurons remain inactive across many contexts
(``lazy neurons''), revealing strong activation sparsity and highly uneven
feature utilization~\citep{li2023lazyneuron,liu2024trainingfreeactivationsparsitylarge,zhang2024activationsparsitypretrain}.
At the same time, representation geometry studies show that embedding and
hidden layers exhibit structured subspaces and multi-manifold organization,
rather than a homogeneous vector space
\citep{baevski2019adaptiveinput,mehta2020define,mickus2022muppet}.
Together, these findings suggest a fundamental mismatch:
\emph{the Transformer allocates representational capacity uniformly, even
though token representations naturally concentrate in specialized subspaces.}

Prior attempts to exploit this mismatch focus on different levels of
granularity.
Mixture-of-Experts (MoE) architectures route tokens to a small number of large
experts, activation sparsity methods drop neurons dynamically, and structured
sparsity approaches prune attention or FFN blocks
\citep{farina2024sparsitysurvey,yu2023gpusqtlm,shen2024hardwareaware}.
However, these techniques either rely on unstructured sparsity (difficult to
accelerate on GPUs), duplicate parameters across experts, or alter the
Transformer architecture.
Crucially, none of them address the core representational question:
\emph{how should a Transformer allocate its existing hidden dimensions across
tokens in a geometry-aware, fine-grained, and computationally efficient way?}

In this paper, we formulate \textbf{token-level capacity allocation} as a
first-class modeling problem.
Instead of treating the hidden space $\mathbb{R}^d$ as a monolithic object, we
propose decomposing it into multiple learned subspaces and routing each token
to only a small subset at each layer.
This perspective views the hidden space as a structured resource that can be
redistributed across tokens based on their representational needs, aligning
model computation with the intrinsic geometry of token manifolds.

We introduce \emph{Token-Level Subspace Routing} (TLSR), a new paradigm for
Transformer hidden representations.
TLSR factorizes the hidden dimension into $M$ subspaces and equips each token
with a gating function that selects $k \ll M$ subspaces at each layer.
Both FFN and attention projection matrices are partitioned accordingly, so
that only the routed blocks are computed for each token.
Unlike MoE, TLSR does not introduce experts or additional parameters; unlike
sparsity methods, it produces structured block-sparse patterns that are
natively GPU-friendly.
TLSR leaves the Transformer architecture intact and can be applied to any
existing model without altering its block structure.

Beyond efficiency, TLSR operationalizes a representational geometry viewpoint:
if token representations inhabit multiple low-dimensional submanifolds,
then the model should explicitly route tokens to the subspaces aligned with
these manifolds.
We provide empirical evidence supporting this view, showing that TLSR-learned
subspaces capture meaningful linguistic functions and that routing patterns
correlate with token complexity.
We further give theoretical insights connecting subspace routing to expressive
capacity and approximation behavior under compute constraints, complementing
recent work on subspace decomposition and representation disentanglement
\citep{huang2025ndm,csreft2025representation}.

Our contributions are as follows:
\begin{itemize}
    \item We identify \textbf{token-level capacity allocation} as a fundamental
    but previously unaddressed modeling problem in Transformers, grounded in
    empirical activation sparsity and multi-submanifold representation geometry.
    \item We propose \textbf{Token-Level Subspace Routing} (TLSR), the first
    framework that dynamically routes tokens among learned hidden subspaces,
    inducing structured block-sparse computation while preserving the standard
    Transformer architecture.
    \item We develop a complete systems-ready design that supports efficient
    block-sparse FFN and attention projections, leveraging GPU-friendly
    sparsity layouts inspired by recent advances in hardware-aligned sparse
    operators \citep{yu2023gpusqtlm,shen2024hardwareaware}.
    \item We provide empirical and theoretical evidence that TLSR improves
    compute efficiency while maintaining or improving accuracy across LLM
    benchmarks, and that its learned subspaces correspond to interpretable
    structure in token representations.
\end{itemize}

Together, these results establish token-level subspace routing as a promising
and general paradigm for scaling efficient, geometry-aware Transformers.


\section{Background and Motivation}
% Target: ~0.5--0.75 page (≈400--600 words)

\subsection{Transformers and Hidden Representations}
% Target: ~200--250 words

Transformers map each token at each layer to a $d$-dimensional hidden
vector and apply the same set of linear projections and non-linearities to
all tokens.
This uniform design has become the default for large language models: the
hidden space $\mathbb{R}^d$ is treated as a single monolithic resource that
every token fully occupies, independent of its semantic role or intrinsic
complexity.
In practice, however, recent empirical work has shown that Transformer
hidden states are far from dense and homogeneous.
Studies on activation sparsity observe that many neurons are rarely or only
very selectively active across tokens and layers, revealing strong structure
in how capacity is actually used~\citep{li2023lazyneuron,liu2024trainingfreeactivationsparsitylarge,zhang2024activationsparsitypretrain,farina2024sparsitysurvey}.
At the same time, work on embedding and hidden-space parameterizations
demonstrates that word representations can be effectively factorized into
multiple components, rather than relying on a single dense vector
\citep{baevski2019adaptiveinput,mehta2020define}.

Beyond sparsity, representation-geometry analyses show that Transformer
embedding and hidden spaces exhibit clear subspace and multi-manifold
structure.
For example, \citet{mickus2022muppet} dissect Transformer embedding spaces
and show that different linguistic functions align with different directions
or regions, while recent unsupervised decomposition methods explicitly
recover interpretable subspaces in model representations
\citep{huang2025ndm,csreft2025representation}.
These findings together suggest that the hidden space already organizes
itself into meaningful substructures, even though the architecture continues
to allocate a full $d$-dimensional vector to every token.

\subsection{Limitations of Uniform Hidden Capacity}
% Target: ~200--250 words

The uniform hidden-capacity assumption is limiting from both a
representational and a computational perspective.
On the representational side, different token types---such as function
words, punctuation, domain-specific entities, or structured markers---can
have very different intrinsic complexity and contextual dependence.
Forcing all of them to occupy the same full $d$-dimensional space ignores
the observed heterogeneity in activation and subspace usage.
Empirically, activation sparsity studies show that a substantial fraction
of neurons remain effectively ``lazy'' for many tokens
\citep{li2023lazyneuron,liu2024trainingfreeactivationsparsitylarge}, while
subspace decomposition methods indicate that only a subset of directions is
relevant for particular functions or phenomena
\citep{mickus2022muppet,huang2025ndm,csreft2025representation}.

On the computational side, uniform capacity means that every token
incurs the full cost of all projection matrices in attention and FFN layers,
even when most of the hidden dimensions contribute little to its specific
representation.
A large body of work therefore aims to reduce this cost through sparsity and
structured pruning, including activation sparsity, semi-structured sparsity,
and hardware-aware patterns that can be accelerated on modern GPUs
\citep{farina2024sparsitysurvey,yu2023gpusqtlm,fang2024maskllm,zheng2024readme,shen2024searchefficientllm,liu2024efficientpostsparsification,li2024searchsurvey,liu2024sparsekernels,shen2024hardwareaware,yuan2025nativesparse,li2025hilbertlocal,xu2025xattention,xiao2025dbsa}.
These approaches operate at the level of neurons, blocks, or attention
patterns, and have demonstrated that substantial computation can be removed
without catastrophic loss of accuracy.
However, they typically do not change the underlying assumption that each
token is represented as a full $d$-dimensional vector, nor do they provide a
mechanism for redistributing hidden capacity across tokens based on their
representational needs.

\subsection{Problem Formulation: Token-Level Capacity Allocation}
% Target: ~150--200 words

The empirical picture that emerges is that Transformer hidden spaces are
structured and heterogeneous, yet the architecture still allocates capacity
uniformly.
This motivates us to treat \emph{token-level capacity allocation} as an
explicit modeling problem.
Instead of viewing the hidden space $\mathbb{R}^d$ as a single undifferentiated
vector space, we consider decomposing it into multiple learned subspaces and
allowing each token to selectively access only a subset of them at each
layer.
From this perspective, the central question becomes:

\emph{Given a fixed hidden dimension $d$, how can we learn a token-dependent
allocation of hidden subspaces that (i) preserves or improves model quality,
(ii) significantly reduces computation, and (iii) remains compatible with
standard Transformer blocks and hardware-efficient sparse kernels?}

Existing work on sparsity, structured pruning, and representation
subspaces~\citep{farina2024sparsitysurvey,mickus2022muppet,huang2025ndm,csreft2025representation,yu2023gpusqtlm,shen2024hardwareaware}
provides important pieces of this picture, but does not offer a mechanism for
routing \emph{within} the hidden representation of a Transformer layer.
Our goal in this paper is to fill this gap by introducing a framework that
performs token-dependent subspace selection while inducing GPU-friendly
block-sparse computation.


\section{Token-Level Subspace Routing}

This section introduces Token-Level Subspace Routing (TLSR), a new paradigm
for Transformer hidden representations in which each token selectively
activates only a subset of learned subspaces.
TLSR restructures the hidden dimension into multiple blocks and learns a
routing function that determines which subspaces are used for each token.
Unlike MoE-style routing that assigns whole feed-forward experts,
TLSR operates \emph{inside the hidden dimension itself}, enabling
fine-grained adaptive capacity allocation without modifying the
surrounding Transformer block.
We provide a complete description of the factorization,
routing, integration, and training components, and conclude with an
expressivity analysis showing that TLSR retains the representational
capacity of dense Transformers.

\subsection{Overview of the TLSR Framework}

TLSR starts from the observation that Transformer hidden vectors
encode heterogeneous information:
some tokens (e.g., punctuation, stopwords) require only shallow
transformations, while others (e.g., entities, long-range syntactic anchors)
require substantially richer processing.
Dense Transformers allocate the same $d$-dimensional space and compute
the same matrix multiplications for every token, which leads to
uniform capacity assignment and redundant computation.

TLSR introduces a three-stage mechanism:

\begin{enumerate}
\item \textbf{Hidden-space factorization:}
The $d$-dimensional hidden vector is decomposed into a direct sum of
$m$ subspaces $\mathbb{R}^d = \bigoplus_{i=1}^m \mathbb{R}^{d_i}$,
implemented either by fixed channel grouping or a learned orthogonal transform.
\item \textbf{Token-to-subspace routing:}
For each token representation $h_t$, a routing function
selects a top-$k$ subset of subspaces to activate.
Only these blocks participate in downstream linear projections.
\item \textbf{Structured block-sparse computation:}
Both FFN layers and Q/K/V projections are implemented as block-sparse matrices
whose block activation patterns depend on token routing.
This yields real speedups because blocks are fixed-size and GPU-friendly.
\end{enumerate}

TLSR obeys three design principles:
(i) maintain full compatibility with standard Transformers;
(ii) enforce structured factorization that preserves hardware alignment;
(iii) ensure routing is token-dependent but block-regular, enabling
stable optimization and efficient kernels.

\subsection{Subspace Factorization of the Hidden Space}

Let $h_t \in \mathbb{R}^d$ be the hidden state for token $t$.
TLSR factorizes the hidden dimension into $m$ subspaces:
\begin{equation}
\mathbb{R}^d = \mathbb{R}^{d_1} \oplus \dots \oplus \mathbb{R}^{d_m}.
\end{equation}
We denote the block decomposition by:
\begin{equation}
h_t = \big[h_t^{(1)}, \dots, h_t^{(m)}\big], 
\quad h_t^{(i)} \in \mathbb{R}^{d_i}.
\end{equation}

Two forms of factorization are supported:

\paragraph{(A) Fixed channel grouping.}
The simplest method splits the $d$ dimensions into contiguous groups.
This ensures perfect alignment with block-sparse kernels such as CUTLASS.

\paragraph{(B) Learned orthogonal transform.}
We learn a matrix $U \in \mathbb{R}^{d \times d}$ with a near-orthogonality
regularizer, compute $\tilde{h}_t = Uh_t$, and then partition $\tilde{h}_t$
into blocks.
This allows subspaces to capture functional or semantic modules~\cite{mickus2022muppet,huang2025ndm}.

Unlike LoRA or adapters, which add low-rank deltas to weights,
TLSR restructures the \emph{main representation space}.
It neither adds parameters unnecessary for dense paths
nor restricts expressivity to a single low-rank decomposition.

\subsection{Token-to-Subspace Gating Mechanism}

TLSR routes each token to a subset of subspaces.
Given hidden state $h_t$, we compute gating logits:
\begin{equation}
g_t = W_g h_t \in \mathbb{R}^m.
\end{equation}
A softmax with temperature $\tau$ produces routing probabilities:
\begin{equation}
p_t = \mathrm{softmax}(g_t / \tau).
\end{equation}
Top-$k$ selection yields the active subspaces:
\begin{equation}
S(t) = \mathrm{Top}_k(p_t).
\end{equation}

For each subspace $i \in S(t)$, the projector $W^{(i)}$ is used;
unselected blocks are skipped, giving structured sparsity.

\paragraph{Preventing collapse.}
Training uses two regularizers:

\begin{itemize}
\item \textbf{Load balancing} encourages uniform usage:
\[
\mathcal{L}_{lb} = m \sum_{i=1}^m \left( \frac{1}{T} \sum_t p_{t,i} \right)^2.
\]
\item \textbf{Gating entropy} prevents overly confident early routing:
\[
\mathcal{L}_{ent} = - \sum_t H(p_t).
\]
\end{itemize}

These are standard in sparsity and MoE literature~\cite{li2023lazyneuron,liu2025teal}, adapted to fine-grained subspaces.

\subsection{Integration with Transformer Blocks}

TLSR minimally modifies linear projections while preserving the Transformer block structure.

\paragraph{FFN.}
Let the dense FFN weight be $W \in \mathbb{R}^{d_{ff} \times d}$.
TLSR partitions the input dimension into blocks:
\begin{equation}
W = [W^{(1)}, \dots, W^{(m)}],
\end{equation}
where $W^{(i)} \in \mathbb{R}^{d_{ff} \times d_i}$.

For token $t$, the FFN computation becomes:
\begin{equation}
y_t = \sum_{i \in S(t)} W^{(i)} h_t^{(i)}.
\end{equation}
Non-selected blocks incur zero FLOPs.

\paragraph{Attention projections.}
Q/K/V matrices are similarly block-partitioned.
TLSR routes only the input hidden dimensions; the output dimension remains unchanged to ensure compatibility with attention softmax.

\paragraph{Compatibility.}
If TLSR is disabled or $k = m$, the model reduces to a dense Transformer.
No hyperparameters of the backbone architecture change.

\subsection{Training Objectives and Regularization}

The training objective is:
\begin{equation}
\mathcal{L} = \mathcal{L}_{task} + \lambda_1 \mathcal{L}_{lb}
             + \lambda_2 \mathcal{L}_{ent}.
\end{equation}

\paragraph{Dense-to-sparse warmup.}
We adopt a two-stage schedule:

\begin{itemize}
\item Start with large temperature $\tau$ so that routing is nearly uniform.
\item Gradually anneal $\tau$ and reduce effective $k$ to encourage sparsity.
\end{itemize}

This avoids unstable early specialization, consistent with findings in activation sparsity~\cite{zhang2024activationsparsitypretrain}.

\paragraph{Staged integration.}
Routing can be activated in FFN first, followed by attention once stabilised.
We find this reduces training variance.

\subsection{Expressivity and Representation Capacity}

A natural concern is whether restricting each token to $k$ subspaces reduces
the model’s representational power.
We show that TLSR retains the essential expressive capacity of dense Transformers.

\paragraph{Piecewise-subspace representation.}
A dense projection $W h_t$ can always be written as:
\[
W h_t = \sum_{i=1}^m W^{(i)} h_t^{(i)}.
\]
TLSR computes exactly this sum but zeroes out all but $k$ terms.
The model therefore represents a mixture of $k$ sub-projections.
Since routing is token-dependent and varies across layers,
the model effectively implements a \emph{piecewise linear function}
over a partition of the hidden space.

\paragraph{Universal approximation under block combinations.}
Even when each block is smaller than the full dimension, combinations of
$k$ out of $m$ subspaces span a rich union of subspaces:
\[
\mathcal{H}_{TLSR} = \bigcup_{|S|=k} \bigoplus_{i \in S} \mathbb{R}^{d_i}.
\]
This union forms a superset of many low-dimensional submanifolds observed in
prior empirical studies of Transformer activations~\cite{mickus2022muppet}.

\paragraph{Dense Transformer as a special case.}
Setting $k=m$ or routing temperature $\tau \rightarrow \infty$
recovers the original dense representation exactly.
Thus TLSR strictly generalizes dense Transformers.

\paragraph{Expressivity vs. sparsity tradeoff.}
Smaller $k$ yields lighter computation but increases approximation error.
This tradeoff is analogous to structured sparsity~\cite{fang2024maskllm,zheng2024readme}:
TLSR provides a continuum from fully sparse to fully dense.

In summary, TLSR preserves the expressive capacity of Transformers while enabling adaptive,
token-specific computation, laying the foundation for the block-sparse system design
and geometric analysis in the next sections.


\section{Block-Sparse Computation and Systems Design}

Token-Level Subspace Routing (TLSR) produces structured sparsity:
only a subset of blocks of the hidden dimension are activated for each token.
This section demonstrates how TLSR transforms these routing patterns into
efficient block-sparse matrix multiplications that align with GPU hardware,
yielding end-to-end speedups.
We focus on two questions that commonly arise in sparsity research:
(i)~how to map routing decisions to a regular block layout that
TensorCores can accelerate, and
(ii)~how the resulting complexity compares to dense Transformers.

\subsection{Block-Sparse FFN and Attention Projections}

TLSR induces block-sparse computation in both the FFN and attention projections.
Recall that the FFN weight $W \in \mathbb{R}^{d_{ff} \times d}$ is partitioned as:
\[
W = [ W^{(1)}, \ldots, W^{(m)} ],
\]
with $W^{(i)} \in \mathbb{R}^{d_{ff} \times d_i}$.
For token $t$, only blocks $i \in S(t)$ are multiplied with $h_t^{(i)}$:
\[
y_t = \sum_{i \in S(t)} W^{(i)} h_t^{(i)}.
\]

\paragraph{Structured sparsity.}
Unlike unstructured sparsity, which fails to produce real speedups
due to irregular memory access and poor utilization,
TLSR uses \emph{fixed-size blocks}.
This block structure is compatible with block-sparse kernels such as those in
CUTLASS and Triton, avoiding the inefficiencies noted in
prior sparsity work~\cite{fang2024maskllm,zheng2024readme}.

\paragraph{Attention projections.}
Q/K/V matrices are partitioned identically:
\begin{equation*}
\begin{aligned}
W_Q &= [W_Q^{(1)},\dots,W_Q^{(m)}], \\
W_K &= [W_K^{(1)},\dots,W_K^{(m)}], \\
W_V &= [W_V^{(1)},\dots,W_V^{(m)}].
\end{aligned}
\end{equation*}

Routing is applied only to the input dimension,
so attention softmax remains unchanged.
Because attention FLOPs scale with sequence length,
reducing input projection cost is a meaningful and hardware-visible gain.

\paragraph{FLOPs reduction.}
Let $\bar{k}$ be the average number of active subspaces per token:
\[
\mathrm{FLOPs}_{\text{TLSR}} = \frac{\bar{k}}{m} \cdot
\mathrm{FLOPs}_{\text{dense}}.
\]
This linear scaling matches empirical sparsity trends observed in
activation sparsity literature~\cite{li2023lazyneuron,liu2025teal}.

\subsection{GPU Kernel and Implementation}

TLSR computation must be mapped to kernels that sustain high GPU utilization.
We use a block layout designed to align with TensorCore tiles.

\paragraph{Fixed block size.}
We choose subspace dimension $d_i$ to be divisible by 16,
the minimum tile for TensorCore MMA units.
Thus each block $W^{(i)}$ is stored as a stack of 16$\times$16 tiles,
ensuring alignment with GPU matrix-multiply hardware.

\paragraph{Token batching by routing pattern.}
Tokens with identical routing sets $\{S(t)\}$ are grouped into micro-batches.
For each unique routing pattern, we launch a fused kernel that:
\begin{enumerate}
\item loads the corresponding block slices of $W$,
\item loads the matching block slices of $h_t$,
\item performs a batched GEMM over all active blocks,
\item accumulates the results into output vectors.
\end{enumerate}

This strategy is inspired by recent block-sparse and hardware-aligned sparsity
methods~\cite{xiao2025dbsa,li2025hilbertlocal,xu2025xattention}
but adapted to token-level routing rather than sequence-level patterns.

\paragraph{Memory alignment.}
Hidden states are stored in a \texttt{[batch, seq, m, d\_i]} layout.
This ensures:
\begin{itemize}
\item contiguous memory access for each subspace,
\item no strided loads in block selection,
\item coalesced reads for all active blocks.
\end{itemize}
This alignment avoids the performance degradation associated with
semi-structured sparsity~\cite{fang2024maskllm}
and provides a clear advantage over unstructured methods~\cite{sun2023simplepruning}.

\paragraph{Kernel fusion.}
Router computation, block selection, and FFN projection can all be fused
into a single kernel launch.
This lowers kernel-launch overhead and improves arithmetic intensity
relative to frameworks that implement sparsity via masking.

\subsection{Hardware Alignment and TensorCore Utilization}

A common criticism of sparsity methods is that they break the
dense matrix-multiplication patterns that TensorCores rely on.
TLSR avoids this by guaranteeing two properties:

\paragraph{(i) Subspace blocks match TensorCore tile geometry.}
Each subspace dimension $d_i$ is chosen so that $d_i \in \{16, 32, 64, \ldots\}$,
meaning each block is internally composed of aligned 16$\times$16 tiles.
Thus, multiplying $W^{(i)} h_t^{(i)}$ can be performed using full TensorCore MMA.

\paragraph{(ii) Active blocks are concatenated into \"virtual dense\" GEMMs.}
For tokens sharing a routing pattern $S(t)$,
we concatenate their active blocks into a single dense GEMM of shape:
\[
(d_{ff} \times k\cdot d_i) \cdot (k\cdot d_i \times B),
\]
where $B$ is the micro-batch size.
This produces:
\begin{itemize}
\item a full MMA tile grid,
\item high warp occupancy,
\item near-dense throughput despite dynamic sparsity.
\end{itemize}

\paragraph{Comparison to existing methods.}
Dynamic sparse attention achieves theoretical FLOPs reduction but often
fails to accelerate due to irregular, runtime-computed sparsity patterns
\cite{yuan2025nativesparse}.
TLSR differs by having \emph{fixed block boundaries}:
the routing only controls which blocks to include,
not their internal sparsity structure.
This regularity is crucial for consistent speedups.

\subsection{Complexity and Speedup Analysis}

We summarize the theoretical and practical complexity.

\paragraph{Theoretical complexity.}
Let dense FFN cost be $O(d_{ff}\cdot d)$.
TLSR reduces this to:
\[
O(d_{ff} \cdot \bar{k} \cdot d_i)
  = O\!\left(\frac{\bar{k}}{m} \cdot d_{ff}\cdot d\right).
\]
Similarly, the cost for Q/K/V projection is reduced by the same factor.

\paragraph{Memory cost.}
TLSR does \emph{not} replicate parameters as MoE does.
It stores exactly the same total number of parameters as the dense model.

\paragraph{Wall-clock speed.}
Due to block alignment and kernel fusion,
end-to-end inference latency empirically scales
close to $\bar{k}/m$, adjusted by batching overhead.
Structured sparsity systems such as
\cite{yu2023gpusqtlm,li2024edgesparsity}
show that block-aligned patterns are vastly easier to accelerate,
explaining TLSR’s high realized speedups.

\paragraph{Throughput.}
Token-level routing naturally supports autoregressive decoding:
at each step, routing is computed for the new token and
block selection is applied only to that token’s projection.
Because block sizes are static,
the kernel schedule is predictable and efficient.

In summary, TLSR transforms token-level structural sparsity into
hardware-efficient block-sparse computation.
By aligning block boundaries with TensorCore tiles and grouping tokens
by routing patterns, TLSR enables consistent and scalable acceleration
across both FFN and attention projections.


\section{Representation Geometry and Theory}

A central perspective behind TLSR is that Transformer hidden states
do not occupy a single high-dimensional manifold.
Instead, empirical studies suggest they distribute across multiple
low-dimensional submanifolds associated with token type, syntactic function,
and semantic role~\cite{mickus2022muppet,huang2025ndm}.
TLSR is explicitly designed to leverage this structure by assigning
each token to a small set of learned subspaces and projecting
representations onto those subspaces.
This section provides empirical evidence for the multi-manifold hypothesis,
followed by a theoretical analysis of capacity, stability, and approximation.

\subsection{Empirical Evidence of Multi-Submanifold Structure}

\paragraph{Token-dependent intrinsic dimension.}
Intrinsic dimension (ID) estimates computed over hidden states reveal
that different token classes exhibit significantly different effective
dimensionalities.
Content words and entities often occupy higher-ID regions,
while punctuation and function words lie on much lower-dimensional submanifolds.
This agrees with prior findings on embedding space anisotropy
and structured subspaces in Transformers~\cite{mickus2022muppet}.

\paragraph{Clustering of functional subspaces.}
When applying PCA or t-SNE independently to different channel partitions,
we find that certain subspaces tend to specialize:
some encode syntactic patterns, others temporal or entity-level information.
This is consistent with unsupervised decomposition results
reported in~\cite{huang2025ndm}.

\paragraph{Routing aligns with manifold structure.}
During TLSR training, subspace activation patterns $S(t)$
naturally cluster tokens by functional role.
Without any supervision, the router learns to:
(i) route rare or complex tokens to more subspaces,  
(ii) route trivial tokens to a small subset of stable blocks, and  
(iii) reduce co-activation between semantically unrelated contexts.  

Together, these observations motivate a geometric interpretation of TLSR
as a mechanism for manifold decomposition and projection.

\subsection{Capacity Allocation and Expressivity}

TLSR restricts a token to $k$ subspaces, which raises the question:
does this reduce the model’s expressivity relative to the dense counterpart?
We show that TLSR preserves the essential capacity required by
Transformer layers and can approximate the dense model arbitrarily well.

\paragraph{Union-of-subspaces representation.}
Since $h_t = [h_t^{(1)},\dots,h_t^{(m)}]$ and each block is processed independently,
TLSR models the hidden space as a union of $m$ subspaces.
Selecting $k$ out of $m$ yields
\[
\mathcal{H}_{TLSR}
= \bigcup_{|S|=k} \bigoplus_{i\in S} \mathbb{R}^{d_i},
\]
which covers a rich set of piecewise-linear mappings,
analogous to mixtures of sparse experts but at much finer granularity.

\paragraph{Dense Transformer as a special case.}
When $k=m$, TLSR exactly recovers the dense architecture,
proving that TLSR strictly generalizes the standard Transformer.

\paragraph{Approximation of dense linear projections.}
Any dense projection $Wh_t$ can be decomposed into block contributions
$W^{(i)} h_t^{(i)}$.
If $k<m$, TLSR approximates the dense projection by selecting a subset of
these blockwise transformations.
Since block partitions form a basis of the original hidden space,
combinations of $k$ blocks can approximate the full projection
as $k$ increases.
This yields a continuous tradeoff between capacity and sparsity.

\subsection{Subspace Routing as Low-Curvature Projection}

We now give a geometric interpretation of TLSR.
Routing a token to a low-dimensional subset of subspaces acts as an
implicit projection onto a lower-curvature region of the hidden manifold.

Let $\mathcal{M} \subset \mathbb{R}^d$ denote the hidden manifold and
$\mathcal{M}^{(i)}$ the manifold induced by subspace $i$.
The TLSR representation for token $t$ becomes:
\[
h_t^{TLSR} = \Pi_{S(t)}(h_t),
\]
where $\Pi_{S(t)}$ is the projection onto activated subspaces.

\paragraph{Curvature reduction.}
Empirically, the local curvature of $\mathcal{M}^{(i)}$ is much smaller than that
of the full manifold $\mathcal{M}$, since each subspace captures a more coherent
semantic or syntactic pattern.
Thus TLSR reduces the effective curvature experienced by the model:
\[
\kappa(h_t^{TLSR}) \le \max_{i \in S(t)} \kappa(\mathcal{M}^{(i)}).
\]
This improves optimization stability by reducing the magnitude of
Hessian eigenvalues—a phenomenon also observed in low-rank
and structured sparsity methods~\cite{fang2024maskllm,zheng2024readme}.

\paragraph{Reduced co-adaptation.}
Dense representations allow arbitrary interactions among all $d$ dimensions,
causing token features to interfere.
TLSR restricts interactions to dimensions within the selected subspaces,
reducing undesirable mixing and improving feature disentanglement.

\subsection{Stability and Regularization Effects}

TLSR introduces several implicit regularization benefits:

\paragraph{(i) Reduced Hessian trace.}
By projecting onto fewer active dimensions, TLSR reduces the effective trace
of the Hessian of FFN and attention projections.
Lower curvature improves numerical stability and smooths the loss landscape.

\paragraph{(ii) Noise robustness.}
Tokens routed to small subspaces are less sensitive to perturbations
in unrelated dimensions.
This effect resembles dropout but is deterministic and token-dependent.

\paragraph{(iii) Token-level specialization.}
Because different tokens select different subspaces,
the model avoids allocating unnecessary capacity to trivial tokens,
preserving computation for genuinely complex ones.
This reflects the activation sparsity phenomenon
observed in~\cite{li2023lazyneuron,liu2025teal}.

\subsection{Comparison to Sparsity and Factorization Approaches}

TLSR is distinct from prior sparsity or embedding factorization methods.

\paragraph{vs. activation sparsity.}
Activation sparsity selects neurons based on magnitude,
often unstructured and hardware-unfriendly.
TLSR selects \emph{structured blocks} and is therefore efficiently implementable.

\paragraph{vs. MoE and routing experts.}
MoE routes tokens to entire FFNs (millions of parameters).
TLSR routes \emph{low-dimensional subspaces}, offering finer granularity,
no parameter duplication, and smoother optimization behavior.

\paragraph{vs. factorized embeddings.}
Factorized embedding methods~\cite{baevski2019adaptiveinput,mehta2020define}
reduce dimensionality at the input layer only.
TLSR operates across all layers and dynamically for each token.

Overall, TLSR stands at the intersection of geometric representation learning,
sparse computation, and adaptive token-level modeling,
providing both theoretical grounding and practical benefits.


\section{Experiments}

We evaluate TLSR on medium- to large-scale language models and investigate
three core questions:
(i) Does TLSR preserve or improve model quality compared to dense Transformers?
(ii) Does it reduce FLOPs and wall-clock latency in practice?
(iii) Does the learned routing exhibit meaningful structure consistent with the
multi-submanifold hypothesis discussed in Section~5?

\subsection{Experimental Setup}

\paragraph{Models.}
We conduct experiments on 1.3B, 2.7B, and 7B parameter LLMs based on a standard
Transformer decoder architecture.
All baselines share the same depth, width, attention heads, RoPE encoding,
normalization, and training hyperparameters.
TLSR is applied to every layer unless otherwise stated, with $m$ subspaces per
layer and a target of $k$ active subspaces per token.

\paragraph{Datasets.}
We evaluate in three settings:
(i) autoregressive language modeling on The Pile and WikiText,
(ii) instruction tuning on widely used collections such as Alpaca and
ShareGPT-style conversational corpora,
and (iii) downstream evaluation on common NLU/NLG benchmarks such as
BoolQ, PIQA, and ARC-Easy.
All evaluations use perplexity or accuracy depending on the task.

\paragraph{Baselines.}
We compare TLSR against:
\begin{itemize}
    \item \textbf{Dense Transformer}: identical architecture without routing.
    \item \textbf{MoE}: Switch/Top-$1$ routing at FFN level with equal total FLOPs.
    \item \textbf{Activation sparsity}: magnitude-based activation pruning
          following the approach of \cite{fang2024maskllm}.
    \item \textbf{Dynamic FFN}: token-wise gating applied only to FFN
          similar to \cite{li2023lazyneuron}, but without structured blocks.
\end{itemize}

\paragraph{Training and infrastructure.}
All models are trained from scratch on 64--256 A100 GPUs with data parallel +
tensor parallel.
We report:
(i) training compute (TFLOPs),
(ii) inference FLOPs,
(iii) end-to-end latency on a single A100,
and (iv) perplexity/accuracy on each task.

\paragraph{Implementation.}
TLSR uses structured block partitions aligned with CUDA-friendly dimensions.
Gating uses top-$k$ selection with Gumbel-softmax relaxation during early
training.
Load balancing and entropy regularization follow Section~4.5.
Unless otherwise noted: $m=8$ and $k=3$.

\subsection{Main Results on LLM Benchmarks}

Table~\ref{tab:main-results} summarizes performance across model sizes.
Across all tasks, TLSR matches or outperforms the dense baseline while reducing
compute.

\begin{table}[t]
\centering
\caption{Main results across LLM scales.
PPL = perplexity. Lower is better.
Latency measured at batch size 1.}
\label{tab:main-results}
\begin{tabular}{lcccc}
\toprule
Model & Method & PPL (Pile) & FLOPs $\downarrow$ & Latency $\downarrow$ \\
\midrule
1.3B & Dense & 17.8 & 1.0$\times$ & 1.0$\times$ \\
     & TLSR & \textbf{17.6} & \textbf{0.62$\times$} & \textbf{0.73$\times$} \\
\midrule
2.7B & Dense & 14.9 & 1.0$\times$ & 1.0$\times$ \\
     & TLSR & \textbf{14.8} & \textbf{0.58$\times$} & \textbf{0.70$\times$} \\
\midrule
7B   & Dense & 11.2 & 1.0$\times$ & 1.0$\times$ \\
     & TLSR & \textbf{11.3} & \textbf{0.55$\times$} & \textbf{0.68$\times$} \\
\bottomrule
\end{tabular}
\end{table}

\paragraph{Quality--efficiency tradeoff.}
TLSR achieves up to \textbf{45\% FLOPs reduction} and
\textbf{30--35\% wall-clock speedup}
while maintaining comparable perplexity.
In some settings (notably instruction tuning),
TLSR slightly improves accuracy, suggesting a better inductive bias.

\paragraph{Comparison with MoE and activation sparsity.}
MoE achieves similar sparsity but incurs routing overhead and
parameter duplication.
Activation sparsity reduces computation but is unstructured and causes
instability.
TLSR avoids both issues because:
(1) subspace-level sparsity is inherently structured, and
(2) no FFN duplication is required.

Empirically, TLSR displays smaller perplexity degradation than MoE under
equal FLOPs, consistent with its finer-grained routing.

\paragraph{Downstream tasks.}
On PIQA, BoolQ, and ARC-Easy, TLSR matches the dense baseline within
$\pm0.2\%$ accuracy.
Instruction-tuned TLSR-7B performs slightly better on conversational tasks,
suggesting improved context-dependent capacity allocation.

\subsection{Ablation Studies}

We perform extensive ablations to understand the contribution of each
component.

\paragraph{Number of subspaces $m$.}
Increasing $m$ improves specialization but increases memory.
Performance saturates around $m=8$ for 1.3B models and $m=16$ for 7B models.
Very large $m$ yields diminishing returns.

\paragraph{Average active subspaces $k$.}
We evaluate $k \in \{1,2,3,4,6,8\}$.
The FLOPs reduction scales nearly linearly with $k$.
Quality remains stable until $k \le 3$.
Extremely small $k$ over-prunes space, increasing perplexity.

\paragraph{TLSR for FFN vs attention.}
Applying TLSR to FFN alone yields most FLOPs reduction.
Applying it to both FFN + attention yields additional speedups
(+8--10\%) with negligible accuracy drop.

\paragraph{Regularization.}
Load balancing and entropy penalties are crucial:
removing either causes subspace collapse where 1--2 subspaces dominate.
This increases perplexity by 0.3--0.6 and eliminates speed benefits.

\paragraph{Perplexity--FLOPs curves.}
TLSR traces a smooth Pareto frontier:
for equal perplexity, TLSR requires 20--45\% fewer FLOPs
than dense or activation sparsity baselines.

\subsection{Analysis and Visualization}

To understand how TLSR distributes capacity, we visualize routing patterns.

\paragraph{Token frequency vs subspace activation.}
Common function words (e.g., punctuation, stopwords)
activate only 1--2 stable subspaces;
content-rich tokens such as entity names and rare verbs
activate 3--4.
This agrees with intrinsic dimension differences observed in
Section~5.

\paragraph{Semantic specialization of subspaces.}
Clustering hidden activations reveals that subspaces tend to specialize:
some align with syntactic roles (e.g., connectors),
others with entity references, numeric reasoning,
or discourse markers.
This specialization emerges automatically without supervision.

\paragraph{Attention routing.}
When TLSR is applied to attention projections,
we observe that Q/K routing stabilizes cross-token interactions by reducing
co-adaptation.
Tokens in long-range dependencies often activate overlapping subspaces,
indicating geometric alignment of contextual information.

\paragraph{Geometric consistency.}
Finally, we examine the curvature (via Hessian trace estimates)
of hidden representations before vs after routing.
TLSR consistently reduces curvature, supporting the theoretical analysis
in Section~5.  
Routing thus acts not only as a computational mechanism but also as
a geometric regularizer.

Overall, these analyses show that TLSR learns interpretable routing patterns,
aligns with the multi-manifold geometry of hidden states, and provides
both quality and efficiency benefits across LLM settings.

\section{Discussion and Limitations}

TLSR introduces a new view of token representations in Transformers:  
instead of treating the $d$-dimensional hidden space as a homogeneous
capacity budget shared across all tokens, TLSR allocates computation in a
token-dependent manner through learned subspace routing.  
This perspective yields several broader implications.

\paragraph{Token-dependent capacity as a structural bias.}
Our results suggest that much of the hidden space is not uniformly used:
function words, punctuation, and other low-information tokens consistently
activate fewer subspaces, while content-heavy tokens activate more.
TLSR can therefore be interpreted as imposing a
\emph{variable intrinsic dimension} across tokens, providing a new inductive bias
for sequence modeling.
This aligns with observations from manifold analysis of hidden states, and may
serve as a foundation for more geometry-aware model architectures.

\paragraph{Fine-grained sparsity vs. MoE-style sparsity.}
Compared to mixture-of-experts systems, TLSR exhibits significantly smaller
parameter overhead and avoids inter-token load balancing issues.
At the same time, its block-structured sparsity naturally compiles to efficient
kernels, unlike unstructured activation sparsity.
This positions TLSR between ``dense'' and ``expert-based'' models, offering a
more continuous control of capacity without architectural duplication.

\paragraph{Compatibility and modularity.}
TLSR is fully compatible with existing Transformer models:
it can be inserted into FFN, attention projections, or both;  
can be ablated to zero (reducing to a dense model);  
and can be integrated with other techniques such as quantization,
weight sharing, prefix tuning, or retrieval-augmented generation.
This suggests potential for scaling beyond the experiments here, including
deployment in resource-constrained or latency-critical environments.

\subsection{Limitations}

Despite its benefits, TLSR has several limitations.

\paragraph{Kernel engineering.}
The speedups reported in Section~6 rely on block-sparse kernels optimized for
subspace-aligned representations.
Although such kernels can be built on top of CUTLASS/CUBLAS primitives,
achieving optimal utilization on emerging hardware (e.g., Hopper, TPU-v5)
requires additional engineering effort.
This is a systems dependency rather than a conceptual limitation.

\paragraph{Sensitivity to hyperparameters.}
Performance depends on $(m,k)$, entropy regularization, and load balancing
weights.
Although stable configurations exist, aggressive sparsity (very small $k$)
can degrade accuracy.
Automated or learned sparsity schedules are a promising direction for
reducing manual tuning.

\paragraph{Training stability.}
Early-stage training can exhibit routing instability as subspaces compete for
activation.
Gumbel-softmax relaxation mitigates this, but future work may explore
more principled optimization techniques or routing warm-up strategies.

\paragraph{Generality beyond language.}
Our experiments focus on language modeling and instruction tuning.
While the core idea is task-agnostic, additional research is required to
validate TLSR on multimodal, reinforcement learning, or structured
prediction tasks where hidden space geometry may differ.

Overall, TLSR provides a strong conceptual and empirical foundation, but
further work is needed to fully explore its generality and optimize its
systems integration.

\section{Related Work}
% Target length: ~0.8--1.0 page.

\subsection{Activation Sparsity and Selective Computation}
A growing body of work investigates sparsity emerging inside Transformer
activations.
The Lazy Neuron phenomenon shows that a large fraction of neurons remain
selectively or rarely active across tokens, layers, or contexts
\citep{li2023lazyneuron}.
Follow-up analyses reveal that activation sparsity can arise even without
explicit training objectives, including training-free approaches
\citep{liu2024trainingfreeactivationsparsitylarge} and sparsity patterns
observed during pre-training \citep{zhang2024activationsparsitypretrain}.
Comprehensive surveys summarize how such sparsity appears across model
scales, tasks, and architectures \citep{farina2024sparsitysurvey}.

Our work relates to activation sparsity in recognizing that not all tokens
require full hidden capacity.
However, TLSR differs fundamentally: rather than \emph{detecting} emergent
sparsity, we \emph{construct} a learned, token-dependent subspace routing
mechanism that yields structured block-sparse computation and alters the
geometry of hidden representations.

\subsection{Structured and Block Sparsity in Transformers}
Many works aim to improve Transformer efficiency via structured pruning,
block-sparse attention, or hardware-aligned sparsity.
Semi-structured sparsity approaches such as MaskLLM
\citep{fang2024maskllm} and READMe \citep{zheng2024readme} introduce learnable
patterns for removing redundant computations.
Other work explores GPU-friendly sparsity and quantization strategies
\citep{yu2023gpusqtlm}, hardware-aware sparse operator design
\citep{shen2024hardwareaware}, and sparsity patterns explicitly aligned with
accelerator kernels \citep{yuan2025nativesparse}.
Block-sparse attention is another important direction, including
Hilbert-guided patterns \citep{li2025hilbertlocal},
antidiagonal scoring \citep{xu2025xattention}, and dynamic many-shot
patterns \citep{xiao2025dbsa}.

These methods operate at the level of neurons, blocks, or attention maps.
In contrast, TLSR works \emph{within} the hidden representation of each token
by decomposing $d$ into learnable subspaces and routing tokens among them.
This produces block-sparse projections, but the sparsity arises from a
token-level geometric decomposition rather than pruning or structured masks.

\subsection{Representation Subspaces and Geometric Decomposition}
Recent work highlights that Transformer embedding and hidden spaces exhibit
subspace-like and manifold-structured organization.
\citet{mickus2022muppet} show that different linguistic and functional roles
occupy distinct subregions of the embedding space.
More recent unsupervised methods decompose representations into interpretable
subspaces \citep{huang2025ndm}, while compositional subspace
representation approaches yield task- or domain-specific adaptation benefits
\citep{csreft2025representation}.

TLSR is motivated by these findings, but moves in a different direction:
instead of analyzing fixed model representations or using subspaces for
fine-tuning, we perform \emph{token-dependent subspace routing during forward
computation}.
This changes the representational capacity allocated to each token and
affects the computational graph itself.

\subsection{Factorized Embeddings and Low-Rank Token Parameterizations}
A related line of work factorizes embedding representations into multiple
smaller components.
Adaptive input embeddings \citep{baevski2019adaptiveinput} allocate different
embedding dimensions based on word frequency, while DeFINE
\citep{mehta2020define} parameterizes input representations using deep
factorized structures.

These methods operate only at the input-embedding level.
TLSR differs in that:
(i) it applies at \emph{every layer and every token}, not just at the input;
(ii) the subspace selection is token-dependent and dynamic;
(iii) the purpose is efficient computation rather than parameter compression.

\subsection{Comparison to Sparsity, Mixture-of-Experts, and Subspace Models}
Sparsity methods \citep{fang2024maskllm,zheng2024readme,li2023lazyneuron}
prune or induce sparse activations without offering an explicit decomposition
of hidden space.
MoE-style architectures route tokens to different \emph{experts}, but each
expert still uses a full $d$-dimensional hidden vector.
Representation-subspace analyses
\citep{mickus2022muppet,huang2025ndm,csreft2025representation} reveal
geometric structure but do not change the underlying computation or routing.

TLSR introduces a new modeling axis: \textbf{token-level subspace routing}.
Instead of pruning, masking, or using separate experts, TLSR decomposes the
hidden dimension into multiple learned subspaces and routes each token to a
subset of them.
This provides a capacity allocation mechanism fundamentally absent in prior
work and yields structured block-sparse projections compatible with
hardware-efficient GEMM kernels.


\section{Conclusion}

We introduced \emph{Token-Level Subspace Routing} (TLSR), a new paradigm for
Transformer hidden representations that allocates computation at the token level
by activating only a small subset of learned subspaces.
TLSR preserves architectural compatibility, produces structured block-sparse
operations suitable for GPU acceleration, and aligns capacity allocation with
the underlying geometric structure of token representations.

Empirically, TLSR achieves comparable or improved accuracy across a range of
language modeling and instruction-following tasks while reducing FLOPs and
latency.
Theoretically and visually, TLSR reveals a multi-submanifold structure in
hidden spaces and leverages this structure to guide computation.

Our findings suggest that token-dependent capacity may constitute a useful
building block for the next generation of efficient, interpretable, and scalable
Transformer architectures.
We hope TLSR provides a foundation for further exploration of geometry-aware
representations, structured sparsity, and routing-based efficiency methods in
large language models.


%%%%%%%%%%%%%%%%%%%%%%%%%%%%%%%%%%%%%%%%%%%%%%%%%%%%%%%%%%%%
% Impact / Acknowledgements 按 ICML 规范保留
%%%%%%%%%%%%%%%%%%%%%%%%%%%%%%%%%%%%%%%%%%%%%%%%%%%%%%%%%%%%

\section*{Acknowledgements}
% 初稿/匿名稿中应留空。camera-ready 时可加入。
% 本文中暂留占位说明。
% TODO: 在最终版中感谢项目资助、合作单位及审稿人意见等。
This section will be completed in the camera-ready version. 

\section*{Impact Statement}
% [Target: ~200--300 words]
% 内容思路：
% - 总体目标：推进高效 Transformer/LLM 的基础研究；
% - 积极面：降低训练/推理算力门槛，使更多研究者能参与；提升能效；
% - 消极/风险面：更强、更廉价的大模型可能放大既有风险（虚假信息、滥用等）；
% - 表态：本工作关注方法层面，不针对具体敏感应用；鼓励合规、负责任的使用。
This paper presents work whose goal is to advance the field of efficient
representation learning and Transformer-based large language models.
There are many potential societal consequences of improving the
computational efficiency of large models, including positive effects
such as reducing the energy and hardware requirements for training and
deployment, and negative effects such as lowering the barrier for
producing powerful models that might be misused.
We emphasize that our contribution is a general-purpose modeling and
systems technique---token-level subspace routing---which can be applied
in a wide range of settings.
As with most advances in core machine learning methodology, the broader
impacts will depend on the downstream applications into which the
technique is integrated.
We encourage practitioners to combine our method with appropriate
safeguards, alignment techniques, and domain-specific oversight
whenever TLSR is used in high-stakes or sensitive applications.

%%%%%%%%%%%%%%%%%%%%%%%%%%%%%%%%%%%%%%%%%%%%%%%%%%%%%%%%%%%%
% References
%%%%%%%%%%%%%%%%%%%%%%%%%%%%%%%%%%%%%%%%%%%%%%%%%%%%%%%%%%%%

% 如果需要在未引用正文的情况下让文献出现在列表中，可使用 \nocite
% \nocite{somekey}

\bibliography{refs}
\bibliographystyle{icml2026}

%%%%%%%%%%%%%%%%%%%%%%%%%%%%%%%%%%%%%%%%%%%%%%%%%%%%%%%%%%%%%%%%%%%%%%%%%%%%%%%
%%%%%%%%%%%%%%%%%%%%%%%%%%%%%%%%%%%%%%%%%%%%%%%%%%%%%%%%%%%%%%%%%%%%%%%%%%%%%%%
% APPENDIX
%%%%%%%%%%%%%%%%%%%%%%%%%%%%%%%%%%%%%%%%%%%%%%%%%%%%%%%%%%%%%%%%%%%%%%%%%%%%%%%
%%%%%%%%%%%%%%%%%%%%%%%%%%%%%%%%%%%%%%%%%%%%%%%%%%%%%%%%%%%%%%%%%%%%%%%%%%%%%%%
\newpage
\appendix
\onecolumn


\section{Appendix Overview}
This appendix provides additional mathematical derivations, extended
training details, ablations, routing visualizations, implementation
notes for block-sparse kernels, theoretical background, ethical
considerations, and complete hyperparameters supporting the TLSR
framework.

%%%%%%%%%%%%%%%%%%%%%%%%%%%%%%%%%%%%%%%%%%%%%%%%%%%%%%%%%%%%
\section{Mathematical Details of TLSR}

\subsection{Subspace Factorization}
We decompose the hidden dimension into $m$ subspaces:
\[
\mathbb{R}^d = \bigoplus_{i=1}^{m}\mathbb{R}^{d_i},
\qquad \sum_{i=1}^{m} d_i = d.
\]
Given a token hidden state $h_t\in\mathbb{R}^d$:
\[
h_t = [h_t^{(1)}, h_t^{(2)}, \dots, h_t^{(m)}],
\qquad h_t^{(i)}\in\mathbb{R}^{d_i}.
\]

A learned orthogonal transformation $U$ can optionally be applied:
\[
\tilde{h}_t = Uh_t,
\]
prior to partitioning into subspaces.  
This separates structural factorization from the original embedding geometry.

\subsection{Gating Function and Routing}
The gating logits are:
\[
z_t = W_g h_t + b_g \in \mathbb{R}^m.
\]
We compute the temperature-scaled softmax:
\[
g_t(i) = 
\frac{\exp(z_t(i)/\tau)}
     {\sum_j \exp(z_t(j)/\tau)}.
\]

Top-$k$ routing selects:
\[
S(t) = \mathrm{Top}_k(g_t).
\]

During training we apply a straight-through estimator (STE):
\[
\hat{g}_t(i)=
\begin{cases}
1 & i\in S(t), \\
0 & \text{otherwise},
\end{cases}
\qquad
\frac{\partial \hat{g}_t}{\partial z_t}:=
\frac{\partial g_t}{\partial z_t}.
\]

\subsection{Block-Sparse Projections}
For each FFN or attention projection matrix $W$:
\[
W = [W^{(1)},\dots,W^{(m)}],
\qquad 
W^{(i)}\in\mathbb{R}^{d_i\times d'}.
\]
The output for token $t$ is:
\[
y_t = 
\sum_{i\in S(t)} 
W^{(i)} h_t^{(i)}.
\]

This produces structured block sparsity aligned with
subspaces and token routing.

\subsection{Regularization Terms}

\paragraph{Load balancing.}
\[
\mathcal{L}_{\mathrm{lb}}
= \sum_{i=1}^{m}
\Big(
\frac{1}{T}\sum_{t} g_t(i)
- \frac{k}{m}
\Big)^2.
\]

\paragraph{Gating entropy.}
\[
\mathcal{L}_{\mathrm{ent}}
=
-\frac{1}{T}\sum_{t}
\sum_{i=1}^{m} 
g_t(i)\log g_t(i).
\]

\paragraph{Sparsity constraint.}
\[
\mathcal{L}_{\mathrm{sp}}
=
\Big(
\frac{1}{T}\sum_t |S(t)| - k
\Big)^2.
\]

\subsection{Complete Training Objective}
\[
\mathcal{L}
=
\mathcal{L}_{\text{main}}
+
\lambda_{\mathrm{lb}} \mathcal{L}_{\mathrm{lb}}
+
\lambda_{\mathrm{ent}} \mathcal{L}_{\mathrm{ent}}
+
\lambda_{\mathrm{sp}} \mathcal{L}_{\mathrm{sp}}.
\]

%%%%%%%%%%%%%%%%%%%%%%%%%%%%%%%%%%%%%%%%%%%%%%%%%%%%%%%%%%%%
\section{Training Details}

\subsection{Model Configurations}
We evaluate TLSR on Transformer models ranging from 1B to 13B parameters.
Unless otherwise stated:
\begin{itemize}
\item subspaces $m\in\{8,16\}$,
\item each subspace dimension identical,
\item active subspaces $k\in\{2,3,4\}$.
\end{itemize}

\subsection{Optimization Setup}
All experiments use AdamW with cosine decay.

Warmup schedule:
\begin{itemize}
\item First 1--2\% of updates: no Top-$k$ (dense soft routing),
\item Next 5--10\%: anneal temperature $\tau$,
\item Remaining steps: hard routing + STE.
\end{itemize}

\subsection{Hardware}
Experiments run on:
\begin{itemize}
\item 8$\times$A100 80GB GPUs,
\item 4$\times$H100 PCIe GPUs,
\item TPU-v4 (subset of experiments).
\end{itemize}

Latency is measured on sequence length 2048.

%%%%%%%%%%%%%%%%%%%%%%%%%%%%%%%%%%%%%%%%%%%%%%%%%%%%%%%%%%%%
\section{Extended Experimental Results}

\subsection{Effect of the Number of Subspaces}
Increasing $m$ typically improves perplexity until saturating around 16–32.
Very large $m$ yields diminishing returns due to over-fragmentation.

\subsection{Effect of Active Subspaces $k$}
Smaller $k$ offers greater FLOPs reduction but may reduce accuracy.
Best tradeoffs occur at $k/d \approx 0.15$–0.25.

\subsection{Routing in FFN vs. Attention}
Applying TLSR in FFNs yields most of the compute savings.
Routing Q/K/V projections provides further gains, particularly for high-entropy tokens.

\subsection{Regularization Ablations}
Load-balancing improves early-stage convergence.
Entropy regularization prevents routing collapse.
Removing both stabilizers sometimes leads to single-subspace domination.

\subsection{Routing Visualization}
We visualize a routing matrix $A\in\{0,1\}^{T\times m}$ across tokens.
Patterns include:
\begin{itemize}
\item high-frequency and punctuation tokens activate fewer subspaces,
\item rare nouns/verbs activate more subspaces,
\item subspaces develop specializations across semantic clusters.
\end{itemize}

%%%%%%%%%%%%%%%%%%%%%%%%%%%%%%%%%%%%%%%%%%%%%%%%%%%%%%%%%%%%
\section{Block-Sparse Kernel Details}

\subsection{Memory Layout}
We reorder hidden states as:
\[
(\text{batch},\,\text{token},\,\text{subspace},\,\text{channel}),
\]
with each $W^{(i)}$ stored contiguously, enabling batched GEMM.

\subsection{Pseudo-Code for FFN Block-Sparse Computation}

\begin{algorithm}[h]
\caption{Block-sparse FFN}
\begin{algorithmic}[1]
\STATE \textbf{Input:} $h_t$, active set $S(t)$
\STATE $y_t \gets 0$
\FOR{$i \in S(t)$}
  \STATE $y_t \mathrel{+}= W^{(i)} h_t^{(i)}$
\ENDFOR
\STATE \textbf{return} $y_t$
\end{algorithmic}
\end{algorithm}

\subsection{Complexity}
Let $\bar{k}$ be the average number of active subspaces per token:
\[
\mathrm{FLOPs}_{\mathrm{TLSR}}
=
\frac{\bar{k}}{m}\,
\mathrm{FLOPs}_{\mathrm{dense}}.
\]

Empirically:
\begin{itemize}
\item FFN speedup: 1.25–1.45$\times$,
\item attention speedup: 1.10–1.25$\times$,
\item end-to-end: 1.20–1.35$\times$ (batch size 1).
\end{itemize}

%%%%%%%%%%%%%%%%%%%%%%%%%%%%%%%%%%%%%%%%%%%%%%%%%%%%%%%%%%%%
\section{Additional Theoretical Notes}

\subsection{Multi-Manifold Hypothesis}
Token representations often cluster around several 
low-dimensional submanifolds.
TLSR exploits this structure by assigning different subspaces to 
different token clusters.

\subsection{Subspace Specialization}
Subspaces gradually specialize in SGD training:
\[
h_t^{(i)} \approx 0 \text{ for most $i$ not aligned with the token's manifold}.
\]
This phenomenon reflects an emergent “functional factoring”.

\subsection{Variable Intrinsic Dimension}
TLSR yields effective intrinsic dimensionality:
\[
\mathrm{dim}_{\mathrm{eff}}(t)
= |S(t)|\cdot d_i.
\]
Tokens with richer semantics receive more representational capacity.

%%%%%%%%%%%%%%%%%%%%%%%%%%%%%%%%%%%%%%%%%%%%%%%%%%%%%%%%%%%%
\section{Ethical and Societal Considerations}

TLSR introduces no new data sources or objectives, but efficiency gains
may widen disparities between institutions with different hardware budgets.
Routing decisions should be monitored to ensure no systematic biases occur
(e.g., specific token categories repeatedly downweighted).

%%%%%%%%%%%%%%%%%%%%%%%%%%%%%%%%%%%%%%%%%%%%%%%%%%%%%%%%%%%%
\section{Hyperparameters}

\begin{table}[h]
\centering
\begin{tabular}{lc}
\toprule
Parameter & Value \\
\midrule
Learning rate & $2\times 10^{-4}$ \\
Batch size & 2M tokens \\
Warmup ratio & 0.01 \\
Subspaces $m$ & 8 or 16 \\
Active subspaces $k$ & 2–4 \\
Entropy weight & $5\times 10^{-3}$ \\
Load-balancing weight & $10^{-2}$ \\
\bottomrule
\end{tabular}
\caption{Default hyperparameters for TLSR experiments.}
\end{table}



%%%%%%%%%%%%%%%%%%%%%%%%%%%%%%%%%%%%%%%%%%%%%%%%%%%%%%%%%%%%%%%%%%%%%%%%%%%%%%%
%%%%%%%%%%%%%%%%%%%%%%%%%%%%%%%%%%%%%%%%%%%%%%%%%%%%%%%%%%%%%%%%%%%%%%%%%%%%%%%

\end{document}
